\section{Dataiku as a workflow reference with explicit scientific acceptance}
\label{sec:exec25-dataiku-reuse}
Dataiku supplies a concrete orchestration reference for the proposed causal
library. The selected evidence concerns the Python API client at revision
\texttt{de83b58cc2d5e1b539fd}\allowbreak\texttt{44797778d56e20f1d367},
committed 3 August 2026, whose package declaration is version 14.7.3. It also
includes retained official DSS 14 and DSS 15 documentation. The immutable
client predates the 7 October 2026 cutoff; the subsequently acquired mutable
documentation is qualified by branch and acquisition date, and does not prove
an exact cutoff server patch or current commercial offer. The reference choice
separates the open client, its dependencies and the commercial server profile.
A client license notice is not a server entitlement, compatibility guarantee
or causal engine qualification.~\citep{exec24_workflow_contracts_dataiku,exec23_dss_dss14_flow,exec23_dss_dss14_releases,exec23_dss_latest_flow,exec23_dss_latest_releases}

This treatment consolidates the preceding Dataiku source findings into an adoption choice; it is not an additional native trial. A useful banking application is a recurring branch-policy analysis. A scenario
can coordinate preparation, fitting, diagnostics and delivery, but a completed
scenario cannot alone establish that the treatment definition, observation
clock, adjustment set and uncertainty calculation are appropriate. Suppose a
monthly analysis is replayed after a data correction. The proposed library
must bind the question version, eligible branch-periods, treatment adoption
record, data and schema digests, method configuration and output identity.
An output from the previous data version is not an accepted answer to the new
request merely because the scenario has a successful execution status. This
scientific acceptance layer is a proposed responsibility of the causal library,
not a claim that the orchestration product lacks a particular undocumented
feature.~\citep{exec25_dataiku_choice}

\subsection{Trigger, run and result are separate objects}
The selected client distinguishes firing a trigger from obtaining a scenario
run. Its \texttt{run} method returns a trigger-fire object; the waiting path
polls for the run or cancellation. Cancellation can return \texttt{None} when
\texttt{no\_fail=True}. In the inspected \texttt{run\_and\_wait} composition,
that result is passed to a waiter without a local \texttt{None} guard, and the
waiter dereferences the scenario-run object. This identifies a conditional
static path to qualify and handle in an adapter; no native cancellation trial
or observed server failure is claimed. The proposed adapter should keep
\emph{trigger cancelled before run creation} distinct from \emph{run aborted}
and from \emph{request completed with a scientific refusal}.~\citep{exec24_workflow_contracts_dataiku}

Completion and outcome also differ. One inspected path regards the presence
of a result as completion, while its strict outcome check accepts
\texttt{SUCCESS} and raises for other outcomes. The documented outcome set
includes \texttt{WARNING}, \texttt{FAILED} and \texttt{ABORTED}; the selected
waiter can return non-success when \texttt{no\_fail} is enabled. Official DSS 14
step-control documentation permits selected failures to be downgraded to
warnings and provides conditional steps, including steps that run regardless
of earlier failure. A proposed causal adapter must retain both the original
outcome and the configured exception policy. A warning permitted by workflow
configuration may still be scientifically disqualifying, while a deliberately
completed refusal can be a valid service response. Neither should be silently
recoded to a successful effect estimate.~\citep{exec24_workflow_contracts_dataiku,exec23_dss_dss14_flow}

The orchestration boundary must therefore expose a logical request identifier,
trigger identity, run identity and attempt identity separately. Consider an
authored example with requests A, B and C: A has two attempts and two created
runs, B has one attempt cancelled before a run, and C has one attempt and one
run. There are three logical requests, four attempts and three runs. The same
run count as request count is accidental; reporting only that count hides both
A's duplicate execution and B's cancellation. No idempotence guarantee is
inferred from the selected client. The proposed replay contract reconciles
existing run identities and accepted artifacts before retrying, and records a
new execution attempt without creating a second scientific request.~\citep{exec25_dataiku_choice}

\subsection{A bounded observer does not prove a stopped computation}
The selected backoff helper starts at 200 milliseconds, multiplies by 1.1 and
caps each sleep at 15 seconds. A cap on an individual polling delay does not
bound the sum of delays. Once the cap is reached, $n$ further waits consume
$15n$ seconds before any HTTP or server time is added. The inspected HTTP call
does not supply a timeout argument, so this source path does not establish a
whole-operation deadline either. These observations are restricted to the
selected source; they are not an assertion about all server or deployment
controls.~\citep{exec24_workflow_contracts_dataiku,exec25_dataiku_choice}

The proposed contract distinguishes a caller's observation deadline from a
confirmed server cancellation. In a stipulated timeline, a trigger fires at
second 0, a run is created at second 3, observation stops at second 10 and the
run completes at second 12. At second 10 the honest state is completion
unknown. An automatic retry could overlap the original run unless the adapter
reconciles its identity and effects. Conversely, a server acknowledgement of
a cancellation request is not accepted as proof that every side effect has
been rolled back. The candidate integration needs bounded client calls,
monotonic total budgets, authenticated terminal-state reconciliation and an
explicit policy for late results. The library must be able to decline a stale
artifact while preserving the costs incurred by the original attempt.~\citep{exec25_dataiku_choice}

\subsection{Conditional integration and the public qualification route}
The present choice is reference use of the selected client and documented
scenario semantics. Conditional integration would wrap a supported and
entitled server/client combination, retain the distinct states above, and
connect completed artifacts to the library's causal question contract. Direct
adoption as a causal engine is not established by orchestration evidence.
Porting selected client logic requires a component and dependency rights
review and reproduces the responsibility for protocol compatibility; porting
proprietary server internals is not authorized by the client notice. Rejection
is specific to unsupported semantic shortcuts, such as equating a successful
scenario with identified causal effect, or treating a local observation timeout
as confirmed remote termination. Other product functionality is outside this
selected-source assessment.~\citep{exec24_workflow_contracts_dataiku,exec25_dataiku_choice}

Before spending on a bank deployment, a public qualification study should
freeze a permitted server profile, compatible client, project and scenario
versions, generated data law, effect query, artifact digests, acceptance rules
and complete cost categories. Generated branch-period fixtures can carry known
causal truth; permitted public datasets can exercise data preparation and
workflow without supplying unobserved assignment truth. The exact scientific
oracle must be supplied by the selected causal method, separately from the
orchestration checks. Fault fixtures include a trigger cancelled before run
creation, a warning-producing step, late completion after an observer deadline,
a stale output digest, duplicate delivery, a failed diagnostic, a scientifically
valid refusal and a changed question version. A fixed direct-library workflow,
the permitted upstream scenario and the proposed adapter should receive the
same scientific task and information. No trial in this protocol has been run
as part of the reference decision.~\citep{exec25_dataiku_choice}

Scientific owners would approve the query, assumptions and acceptance rule;
workflow owners would maintain scenario-to-artifact identity; platform and
procurement owners would establish server rights, compatibility and support;
evaluators would control the fault and causal oracles; maintainers would
review version drift. These remain proposed responsibilities. An illustrative
USD 1,200 preparation cost, USD 1,800 for all attempts, USD 2,400 for scientific
review and USD 600 maintenance allocation total USD 6,000. Across 120 attempted
requests, 90 completions and 60 admissible outputs, the corresponding costs
are USD 50, USD $200/3$ and USD 100. These authored amounts establish accounting
semantics, not a quote or return on investment. The commercial decision requires
actual entitlements, measured complete costs and accepted causal outputs from
the selected supported deployment.~\citep{exec25_dataiku_choice}
